\clearpage
\section{Glossary}

\begin{description}

\item[IGDB]
Internet Game Database; a large structured database of video game information used to programmatically fetch game metadata via its official API.

\item[Launcher]
An interface or application that opens a game executable upon user input.

\item[One-Hot Encoding]
A machine learning technique used to convert categorical data (such as user moods or game genres) into a numerical binary format.

\item[Platforms]
Digital game distribution services (e.g., Steam, GOG, Epic Games Store) where users purchase, subscribe to, and manage games.

\item[TF--IDF Vectorizer]
Term Frequency – Inverse Document Frequency; a technique used to convert text documents (such as game descriptions) into numerical vectors that machine learning models can process.

\end{description}

\section{Analysis Model}

\begin{description}

\item[Data to ML Quantification Pipeline]
The primary analysis model used to process raw system data into machine-readable formats. 
This pipeline converts text descriptions into TF--IDF float vectors, playtime into normalized
values scaled between 0 and 1, and time-of-day information into temporal encodings.

\item[Engagement Score Calculation Model]
An analytical model that computes user engagement with a game by calculating a weighted
combination of the following components:

\begin{itemize}
\item Playtime Score
\item Frequency Score
\item Length Score
\end{itemize}

\item[Hybrid Ranker Model]
The final analysis model responsible for generating recommendations. It calculates a final
recommendation score between 0.0 and 1.0 using the following weighted components:

\begin{itemize}
\item Preference Match (0.40)
\item Similarity to Liked Games (0.35)
\item Quality Score (0.15)
\item Exploration Factor (0.10)
\end{itemize}

\end{description}

\section{Issues List}

\subsection*{Open Issues}

\begin{itemize}

\item The machine learning recommendation module currently yields inconsistent recommendation
quality. The Python implementation remains a work in progress and requires further calibration.

\item The intended Qt/QML user interface design is not yet finalized.

\item The system experiences the \textit{Cold Start} data sparsity problem for new users with zero
recorded playtime. This is currently mitigated by asking users to manually select preferred
genres during the initial setup process.

\end{itemize}

\subsection*{Resolved Issues}

\begin{itemize}

\item The project environment previously failed to place \texttt{std::filesystem} in a namespace.
This issue was resolved by upgrading the project standard from C++14 to C++17.

\item The local heuristic game scanner initially detected redundant executables such as crash
handlers and uninstallers. This issue was resolved by implementing block parameters for
undesired executable names and applying a heuristic that selects the largest executable file.

\end{itemize}
